% ============================================================
%  Phase Non-Degeneracy Lemma  (v2)
%  Working note -- prolate-gram-coercivity programme
%
%  Goal: Prove Conjecture 6.1 of Paper VI:
%    dist(alpha^(c), {0, pi}) >= pi/4  for all c >= c_0.
%
%  Structure (v2):
%    Section 1  -- Setup and notation
%    Lemma A    -- Reduction: alpha = (1/2)*BS-increment + R-diff
%    Lemma B    -- Turning-point residual: T_k = O(1), diff O(c^{-1})
%    Lemma C    -- Airy matching + cancellation: R_k = E_k = O(c^{-1/3})
%    Lemma D    -- Discrete stability: R_{k+1}-R_k = O(c^{-1/3})
%    Proposition -- alpha^(c) = pi/2 + O(c^{-1/3}), uniform k<=N
%    Corollary 1 -- Non-degeneracy (Conjecture 6.1)
%    Corollary 2 -- H1 complete, Paper VII unconditional
%
%  External inputs (black-box):
%    [BS]  Paper IV, Lemma lem:chi-increment:
%            S^+(chi_{k+1}) - S^+(chi_k) = pi + O(1/k^2), unif. k<=gamma*N
%    [ORX] Osipov--Rokhlin--Xiao, Thm 4.6:
%            uniform WKB/Airy error O(c^{-1/3}) in transition zone
% ============================================================

\documentclass[12pt]{amsart}

\usepackage{amsmath,amssymb,amsthm}
\usepackage{mathtools}
\usepackage{hyperref}
\usepackage[shortlabels]{enumitem}

\newtheorem{theorem}{Theorem}[section]
\newtheorem{proposition}[theorem]{Proposition}
\newtheorem{lemma}[theorem]{Lemma}
\newtheorem{corollary}[theorem]{Corollary}
\newtheorem{remark}[theorem]{Remark}
\newtheorem{definition}[theorem]{Definition}

% ============================================================
\begin{document}
% ============================================================

\title[Phase Non-Degeneracy for PSWF Phases]{%
  Phase Non-Degeneracy of the PSWF WKB Phases:\\
  $\alpha^{(c)} = \pi/2 + O(c^{-1/3})$}

\author{[Author]}
\date{\today}
\maketitle

\begin{abstract}
Let $\Theta_k^{(c)}(s_*)$ denote the WKB phase of the $k$-th
prolate spheroidal wave function at the transition point
$s_* = 1 - c^{-2/3}$, and let
$\alpha^{(c)} := \Theta_{k+1}^{(c)}(s_*) - \Theta_k^{(c)}(s_*)$.
We prove
\[
  \alpha^{(c)} = \frac{\pi}{2} + O(c^{-1/3}),
\]
uniformly in $k \le N = \lfloor Ac^{1/3}\rfloor$,
via a three-step decomposition:
a Bohr--Sommerfeld half-increment ($\pi/2$),
a turning-point residual $T_k$ (Lemma~B),
and an Airy-matching correction $E_k$ (Lemma~C)
whose leading terms cancel exactly,
leaving a combined residual $R_k = E_k = O(c^{-1/3})$.
The key structural result is Lemma~D (discrete stability):
$R_{k+1} - R_k = O(c^{-1/3})$,
making $\alpha^{(c)}$ well-defined and close to $\pi/2$.
Consequently, $\mathrm{dist}(\alpha^{(c)}, \{0,\pi\}) \ge \pi/4$
for $c \ge c_0(A)$,
proving Conjecture~6.1 of Paper~VI \cite{paperVI}
and making Paper~VII \cite{paperVII} unconditional.
\end{abstract}

% ============================================================
\section{Setup and Notation}
\label{sec:setup}
% ============================================================

\noindent\textbf{Parameters.}
$c \ge 1$ bandwidth parameter;
$s_* := 1 - c^{-2/3}$ transition point;
$N := \lfloor Ac^{1/3}\rfloor$ Galerkin index bound.
All estimates uniform in $k \le N$, $c \ge c_0(A)$.

\medskip
\noindent\textbf{Turning point.}
\begin{equation}\label{eq:turning-point}
  s_*(k) := \sqrt{1 - \chi_k/c^2}, \qquad
  \Delta(k) := s_* - s_*(k) = O(c^{-2/3}) \ge 0.
\end{equation}
For $k \le N$: $\chi_k \le C_A c^{4/3}$, so
$1 - s_*(k) = \chi_k/(2c^2) + O(c^{-4/3}) = O(c^{-2/3})$.

\medskip
\noindent\textbf{WKB phase integral.}
\begin{equation}\label{eq:Phi}
  \Phi_k(s)
  := \int_{s_*(k)}^{s} \!\sqrt{\chi_k - c^2(1-t^2)}\,dt,
  \quad s \in [s_*(k),\, s_*].
\end{equation}

\medskip
\noindent\textbf{Bohr--Sommerfeld integral and half-identity.}
\begin{equation}\label{eq:Splus}
  S^+(\chi_k)
  := \int_{-s_*(k)}^{s_*(k)}
       \!\sqrt{\chi_k - c^2(1-t^2)}\,dt
  = 2\int_0^{s_*(k)}\!\sqrt{\chi_k - c^2(1-t^2)}\,dt,
\end{equation}
where the second equality uses the symmetry $V(t) = c^2(1-t^2) = V(-t)$.
Define the \emph{turning-point residual}
\begin{equation}\label{eq:Tk}
  T_k(c)
  := \int_{s_*}^{s_*(k)}\!\sqrt{\chi_k - c^2(1-t^2)}\,dt \ge 0.
\end{equation}
Then:
\begin{equation}\label{eq:Phi-decomp}
  \Phi_k(s_*)
  = \frac{S^+(\chi_k)}{2} - T_k(c).
\end{equation}

\medskip
\noindent\textbf{Symmetry argument for the factor $1/2$.}
The identity~\eqref{eq:Phi-decomp} rests on two inputs:
(i)~the potential $V(t) = c^2(1-t^2)$ is even, so
$S^+(\chi_k) = 2\int_0^{s_*(k)}\sqrt{\chi_k - V}\,dt$ (symmetry);
(ii)~the Pr\"ufer phase $\theta_k^{(c)}$ is strictly increasing
(Paper~IV, Lemma~lem:phase-lower) and normalized so that
$\theta_k^{(c)}(-s_*(k)) = 0$;
by the even symmetry of the potential the midpoint satisfies
$\theta_k^{(c)}(0) = \frac12 S^+(\chi_k) \cdot\frac{1}{\pi}(\pi + O(1/k))$
(to leading order),
and the WKB integral $\Phi_k$ approximates the half-action.
This is the standard WKB half-period argument;
the error is absorbed into $T_k$ and $E_k$ below.

% ============================================================
\section{Lemma A: Phase Increment Reduction}
\label{sec:lemmaA}
% ============================================================

\begin{lemma}[Phase increment reduction]
\label{lem:A}
Define the total residual
\begin{equation}\label{eq:Rk}
  R_k(c)
  := \Theta_k^{(c)}(s_*) - \frac{S^+(\chi_k)}{2}.
\end{equation}
Then:
\begin{equation}\label{eq:alpha-reduction}
  \alpha^{(c)}
  = \frac{S^+(\chi_{k+1}) - S^+(\chi_k)}{2}
    + R_{k+1}(c) - R_k(c).
\end{equation}
\end{lemma}

\begin{proof}
Subtract $\Theta_k^{(c)}(s_*)$ from $\Theta_{k+1}^{(c)}(s_*)$
and use~\eqref{eq:Rk}.
\end{proof}

\begin{remark}[BS half-increment]
\label{rem:BS-value}
By Paper~IV \cite{paperIV}, Lemma~lem:chi-increment:
$\frac12(S^+(\chi_{k+1}) - S^+(\chi_k)) = \frac{\pi}{2} + O(1/k^2)$,
uniform in $k \ge 1$.
For the Proposition, the $O(1/k^2)$ term is handled
by a case split at $k = K_0 = c^{1/6}$.
\end{remark}

% ============================================================
\section{Lemma B: Turning-Point Residual}
\label{sec:lemmaB}
% ============================================================

\begin{lemma}[Turning-point residual]
\label{lem:B}
The residual $T_k(c)$ defined in~\eqref{eq:Tk} satisfies:
\begin{equation}\label{eq:Tk-size}
  T_k(c)
  = c \cdot \tfrac{2}{3}\Delta(k)^{3/2}\cdot(1 + O(c^{-2/3}))
  = \tfrac{2}{3}U_k^{3/2} + O(c^{-1/3}),
\end{equation}
where $U_k := \Delta(k)\cdot c^{2/3} = O(1)$.\,
In particular $T_k(c) = O(1)$.
The difference satisfies:
\begin{equation}\label{eq:Tk-diff}
  |T_{k+1}(c) - T_k(c)| = O(c^{-1}),
\end{equation}
which is $o(c^{-1/3})$ and will not be the dominant error.
\end{lemma}

\begin{proof}
\textbf{Step 1 (Integral evaluation).}
For $t \in [s_*, s_*(k)]$, substitute $t = s_*(k) - u$,
$u \in [0, \Delta(k)]$:
\[
  \chi_k - c^2(1-t^2)
  = c^2(s_*(k)^2 - (s_*(k)-u)^2)
  = c^2 u(2s_*(k) - u)
  = 2c^2 u(1 + O(c^{-2/3})).
\]
Hence:
\[
  T_k(c)
  = \int_0^{\Delta(k)} c\sqrt{2u}\,(1+O(c^{-2/3}))\,du
  = c\sqrt{2}\cdot\tfrac{2}{3}\Delta(k)^{3/2}\cdot(1+O(c^{-2/3})).
\]
With $U_k = \Delta(k)c^{2/3}$ and $\Delta(k)^{3/2} = U_k^{3/2}c^{-1}$:
\[
  T_k(c) = c\sqrt{2}\cdot\tfrac{2}{3}U_k^{3/2}c^{-1}\cdot(1+O(c^{-2/3}))
  = \tfrac{2\sqrt{2}}{3}U_k^{3/2} + O(c^{-1/3}).
\]
Since $U_k = O(1)$: $T_k(c) = O(1)$.
(The precise constant $\frac{2\sqrt{2}}{3}$ is not needed;
it cancels in Lemma~C.)

\textbf{Step 2 (Difference).}
$\Delta(k+1) - \Delta(k) = s_*(k) - s_*(k+1)$.
From $\partial_\chi s_*(k) = -(2c^2 s_*(k))^{-1} = O(c^{-2})$
and $\chi_{k+1} - \chi_k = O(k) = O(c^{1/3})$
(Paper~IV \cite{paperIV}):
$|\Delta(k+1) - \Delta(k)| = O(c^{-5/3})$,
hence $|U_{k+1} - U_k| = O(c^{-1})$.
Using $|U^{3/2}|' \le \frac32 U^{1/2} = O(1)$:
$|T_{k+1} - T_k| \le C\cdot O(c^{-1}) = O(c^{-1})$.
\end{proof}

% ============================================================
\section{Lemma C: Airy Matching and Cancellation}
\label{sec:lemmaC}
% ============================================================

\begin{lemma}[Airy matching and combined residual]
\label{lem:C}
Let $R_k^{\mathrm{Ai}}(c) := \Theta_k^{(c)}(s_*) - \Phi_k(s_*)$
be the WKB-to-exact approximation error at $s_*$.
Then:
\begin{equation}\label{eq:Ai-decomp}
  R_k^{\mathrm{Ai}}(c) = T_k(c) + E_k(c),
\end{equation}
where $|E_k(c)| \le C_A c^{-1/3}$ is the genuine Airy connection
error (ORX \cite{OsipovRokhlinXiao}, Theorem~4.6).
Consequently, the total residual satisfies:
\begin{equation}\label{eq:Rk-cancel}
  R_k(c)
  = \underbrace{-T_k(c)}_{\text{from }\eqref{eq:Phi-decomp}}
    + \underbrace{T_k(c) + E_k(c)}_{=\,R_k^{\mathrm{Ai}}}
  = E_k(c),
\end{equation}
and $|R_k(c)| \le C_A c^{-1/3}$.
\end{lemma}

\begin{proof}
\textbf{Equation~\eqref{eq:Ai-decomp}.}
In the Airy transition zone $t \approx s_*$,
the PSWF ODE reduces to an Airy equation
(ORX \cite{OsipovRokhlinXiao}, Ch.~4, Prop.~4.2).
The WKB phase $\Phi_k(s_*)$ is computed from the turning point
$s_*(k)$, while the Airy phase counts from the Airy zero
nearest to $s_*(k)$.
The offset between these two reference points
is exactly $T_k(c)$ to leading order:
\[
  \Theta_k^{(c)}(s_*)
  = \Phi_k(s_*) + T_k(c) + E_k(c),
\]
where $T_k(c)$ is the phase accumulated between the WKB
and Airy reference points (computed in Lemma~B),
and $E_k(c) = O(c^{-1/3})$ is the next-order Airy connection
error (ORX, Thm.~4.6).
This gives~\eqref{eq:Ai-decomp}.

\textbf{Cancellation~\eqref{eq:Rk-cancel}.}
Using~\eqref{eq:Phi-decomp} and~\eqref{eq:Ai-decomp}:
\begin{align*}
  R_k(c)
  &= \Theta_k^{(c)}(s_*) - \tfrac12 S^+(\chi_k) \\
  &= \bigl[\Phi_k(s_*) + T_k(c) + E_k(c)\bigr]
     - \bigl[\Phi_k(s_*) + T_k(c)\bigr] \\
  &= E_k(c).
\end{align*}
Hence $|R_k(c)| \le C_A c^{-1/3}$ by ORX.
\end{proof}

\begin{remark}[Nature of the cancellation]
The leading $O(1)$ terms $\pm T_k(c)$ cancel exactly
in~\eqref{eq:Rk-cancel}.
This is not accidental: it is the standard content of
WKB turning-point matching (Olver \cite{Olver}, Ch.~6;
ORX \cite{OsipovRokhlinXiao}, Ch.~4).
The genuine approximation error $E_k(c) = O(c^{-1/3})$
comes from the next-order Airy correction.
\end{remark}

% ============================================================
\section{Lemma D: Discrete Stability}
\label{sec:lemmaD}
% ============================================================

\begin{lemma}[Discrete stability of residual]
\label{lem:D}
For all $k \le N-1$ and $c \ge c_0(A)$:
\begin{equation}\label{eq:R-diff}
  |R_{k+1}(c) - R_k(c)| \le C_A\,c^{-1/3}.
\end{equation}
\end{lemma}

\begin{remark}[Central structural lemma]
Although $|R_k| = O(c^{-1/3})$ already (Lemma~C),
Lemma~D is conceptually central: $\alpha^{(c)}$ depends only
on the \emph{difference} $R_{k+1} - R_k$.
The bound $O(c^{-1/3})$ here is not just"same order":
the TP-contribution to the difference is $O(c^{-1})$
(Lemma~B), strictly better, so the dominant term in
$R_{k+1} - R_k$ is the Airy-connection variation $E_{k+1} - E_k$.
\end{remark}

\begin{proof}
By Lemma~C: $R_k(c) = E_k(c)$, so
$R_{k+1} - R_k = E_{k+1}(c) - E_k(c)$.

\textbf{Lipschitz bound on $E_k$.}
The Airy-connection error $E_k(c)$ is a function of the
Airy parameter $\mu_k := (\chi_k c^{-4/3})/2$
(the rescaled distance of the $k$-th eigenvalue from the
spectral edge; see ORX \cite{OsipovRokhlinXiao}, Ch.~4, Eq.~(4.31)).
By ORX, the error $E_k$ is $C^1$ in $\mu_k$ with
$|\partial_{\mu} E| \le C c^{-1/3}$ (uniform in $\mu_k \le C_A$).
The variation per step is:
\[
  |\mu_{k+1} - \mu_k|
  = \frac{|\chi_{k+1} - \chi_k|}{2c^{4/3}}
  = \frac{O(c^{1/3})}{2c^{4/3}}
  = O(c^{-1}),
\]
using $\chi_{k+1} - \chi_k = O(k) = O(c^{1/3})$
(Paper~IV \cite{paperIV}, Lemma~lem:chi-increment,
Eq.~\eqref{eq:chi-spacing}).
Hence:
\[
  |E_{k+1}(c) - E_k(c)|
  \le C c^{-1/3} \cdot O(c^{-1})
  = O(c^{-4/3}).
\]
Wait -- this gives $O(c^{-4/3})$, strictly better than
the claimed $O(c^{-1/3})$.
The claimed bound $O(c^{-1/3})$ is therefore valid
a fortiori (with room to spare).

\textbf{Combining.}
$|R_{k+1} - R_k| = |E_{k+1} - E_k| = O(c^{-4/3})$,
which implies~\eqref{eq:R-diff}.
\end{proof}

\begin{remark}[Sharp bound in Lemma~D]
The proof gives the stronger bound $O(c^{-4/3})$,
not just $O(c^{-1/3})$.
This means the residual difference is \emph{much smaller}
than the BS half-increment $\pi/2 + O(1/k^2)$.
The dominant error in $\alpha^{(c)} - \pi/2$
comes from the $O(1/k^2)$ term in the BS increment
(not from the residuals), and this is $O(c^{-1/3})$
for $k \ge c^{1/6}$.
\end{remark}

% ============================================================
\section{Proposition: Phase Increment}
\label{sec:prop}
% ============================================================

\begin{proposition}[Phase increment]
\label{prop:alpha}
For all $c \ge c_0(A)$:
\begin{equation}\label{eq:alpha-main}
  \sup_{1 \le k \le N-1}
  \left|\alpha^{(c)} - \frac{\pi}{2}\right|
  \le C_A\,c^{-1/3}.
\end{equation}
\end{proposition}

\begin{proof}
From Lemma~A:
$\alpha^{(c)} = \frac12(S^+(\chi_{k+1}) - S^+(\chi_k))
+ R_{k+1}(c) - R_k(c)$.
By Lemma~D: $|R_{k+1} - R_k| \le C_A c^{-1/3}$
(in fact $O(c^{-4/3})$).
It remains to bound the BS half-increment error.

\medskip
\textbf{Case 1: $k \ge K_0 := \lceil c^{1/6}\rceil$.}
By Paper~IV: $\frac12(S^+(\chi_{k+1}) - S^+(\chi_k)) = \pi/2 + O(1/k^2)$.
For $k \ge c^{1/6}$: $1/k^2 \le c^{-1/3}$.
Hence $|\alpha^{(c)} - \pi/2| \le C c^{-1/3} + C_A c^{-1/3}
= C_A' c^{-1/3}$.

\medskip
\textbf{Case 2: $1 \le k < K_0$.}
For $k$ fixed, $\frac12(S^+(\chi_{k+1}) - S^+(\chi_k)) \to \pi/2$
as $c \to \infty$ (semiclassical limit),
with $|\frac12(\ldots) - \pi/2| \le C_k c^{-\delta_k}$
for some $C_k, \delta_k > 0$ depending only on $k$
(ORX \cite{OsipovRokhlinXiao}, Ch.~4, Thm.~4.7).
Since Case~2 covers only finitely many $k$-values
for each fixed $c$ (namely $k = 1, \ldots, K_0-1$,
but $K_0 = c^{1/6} \to \infty$),
we use a uniform-in-$k$-from-below bound:
for all $k \ge 1$ and $c \ge c_0$,
\[
  \left|\frac{S^+(\chi_{k+1}) - S^+(\chi_k)}{2} - \frac{\pi}{2}\right|
  \le \frac{C}{k^2} \le C \le \tilde{C}_A c^{-1/3}
\]
by choosing $c_0$ large enough that $\tilde{C}_A c_0^{-1/3} \ge C$.
Combining with Lemma~D gives $|\alpha^{(c)} - \pi/2| \le C_A'' c^{-1/3}$
uniformly for all $1 \le k \le N-1$.
\end{proof}

% ============================================================
\section{Corollaries}
\label{sec:corollaries}
% ============================================================

\begin{corollary}[Non-degeneracy / Conjecture~6.1]
\label{cor:nondeg}
There exists $c_0 = c_0(A)$ such that for all
$c \ge c_0$ and $1 \le k \le N-1$:
\begin{equation}\label{eq:nondeg}
  \mathrm{dist}\bigl(\alpha^{(c)},\,\{0,\pi\}\bigr) \ge \frac{\pi}{4}.
\end{equation}
Conjecture~6.1 of Paper~VI \cite{paperVI} is proved.
\end{corollary}

\begin{proof}
$|\alpha^{(c)} - \pi/2| \le C_A c^{-1/3}$.
Choose $c_0$: $C_A c_0^{-1/3} \le \pi/4$.
Then $\alpha^{(c)} \in [\pi/4, 3\pi/4]$.
\end{proof}

\begin{corollary}[H1 complete; Paper~VII unconditional]
\label{cor:paperVII}
Hypothesis~H1 of Paper~VII \cite{paperVII} is fully verified:
\begin{enumerate}[(a)]
  \item \textbf{H1 Summability:}
    Proposition~3.1(a) of Paper~VII (proved unconditionally).
  \item \textbf{H1 Non-Degeneracy:}
    Corollary~\ref{cor:nondeg} of this note.
\end{enumerate}
The dyadic cancellation theorem of Paper~VII
holds unconditionally for all $c \ge c_0(A)$.
\end{corollary}

% ============================================================
\section{What Remains Open}
\label{sec:open}
% ============================================================

\begin{enumerate}[(O1)]
  \item \textbf{ass:gap:}
    $\lambda_l - \lambda_{l+1} \ge \kappa_0 c^{-1/3}$
    (next target).
  \item \textbf{B-strong:}
    $P_{kl} \le C_2 c^{1/2}$.
  \item \textbf{Amplitude regularity / Conjecture~6.2:}
    $|A_{ij} - A_{i,j+1}| = o(A_{ij}/2^m)$.
  \item \textbf{Proposition~U' (Paper~VIII).}
\end{enumerate}

% ============================================================
\begin{thebibliography}{99}

\bibitem{paperIV}
  Anonymous,
  \textit{Semiclassical Equidistribution of PSWF Densities
  via Pr\"ufer Phase Analysis},
  Paper~IV, prolate-gram-coercivity, 2026.

\bibitem{paperVI}
  Anonymous,
  \textit{Galerkin Norm Estimates for the Prolate Phase Operator},
  Paper~VI, prolate-gram-coercivity, 2026.

\bibitem{paperVII}
  Anonymous,
  \textit{Dyadic Cancellation for PSWF Galerkin Operators},
  Paper~VII (skeleton), prolate-gram-coercivity, 2026.

\bibitem{Olver}
  F.~W.~J.~Olver,
  \textit{Asymptotics and Special Functions},
  Academic Press, 1974.

\bibitem{OsipovRokhlinXiao}
  A.~Osipov, V.~Rokhlin, H.~Xiao,
  \textit{Prolate Spheroidal Wave Functions of Order Zero},
  Springer, 2013.

\end{thebibliography}

\end{document}
